\documentclass[10pt,a4paper]{amsart}
\usepackage[margin=19mm]{geometry}
\usepackage{amsmath,amssymb,mathtools}
\usepackage[scr=rsfs,cal=euler]{mathalfa}
\usepackage{xfrac}
\usepackage[unicode,hidelinks,bookmarks=false]{hyperref}
\newcommand{\ie}{i.e.\ }
\newcommand{\cf}{cf.\ }
\newcommand{\resp}{respectively\ }
\newcommand{\Subsection}{\subsection}
\providecommand{\nobreakdash}{\nobreak\hbox{-}}
\setlength{\emergencystretch}{2em}
\multlinegap0pt
\usepackage{bm}
\usepackage{bbm}
\usepackage{dsfont}
\usepackage{xspace}
\usepackage{cancel}
\usepackage{mathtools}
\usepackage{amsthm}
\makeatletter
\@ifundefined{assumption}{\newtheorem{assumption}{Assumption}}{}
\@ifundefined{lemma}{\newtheorem{lemma}{Lemma}}{}
\@ifundefined{theorem}{\newtheorem{theorem}{Theorem}}{}
\makeatother
\title[Extended HJB Equation for Mean-Variance Stopping Problem: Va]{Extended HJB Equation for Mean-Variance Stopping Problem: Vanishing Regularization Method}
\author{Yuchao Dong, Harry Zheng}
\date{Source: arXiv:2510.24128v1 (CC BY 4.0)}
\begin{document}
\maketitle

\noindent\emph{Source and licence.} Extended HJB Equation for Mean-Variance Stopping Problem: Vanishing Regularization Method. Version arXiv:2510.24128v1 (CC BY 4.0), distributed under the Creative Commons Attribution 4.0 International licence, as recorded in the arXiv licence metadata for this exact version and shown on the abstract page https://arxiv.org/abs/2510.24128v1. Full rights notice: licenses/arxiv-2510.24128-CC-BY-4.0.txt.\par\smallskip

\noindent\textit{Editorial scope.} Counted unit: the Theorem whose source label is \texttt{thm3.2} in the article (source0.tex lines 274--276, source label \texttt{thm3.2}), delivered together with the article's own complete original proof of it (source0.tex lines 277--388, one \texttt{proof} environment of 112 lines). No mathematics is written, repaired or re-derived: statement and proof are the article's own text line for line. Required context retained uncounted: ass (lines 90--101); extended-HJB (lines 173--180); append\_lemma (lines 878--895). Every definition, display and label of those lines is retained unchanged. Omitted from this package and not redistributed: 1--89, 102--172, 181--273, 389--877, 896--938. Nothing inside a retained excerpt is omitted or altered: every retained body is delivered exactly as the article's lines are written, so no omission receipt of any kind accompanies this package. Citations: the retained excerpts carry no citation command at all, so no bibliography entry of the article is transcribed and no reference is redistributed. Reference adaptation: none was needed and none was made; every reference inside the delivered unit resolves inside it, so no unresolved reference or citation is produced. Printed slips kept literal: no printed word, symbol, hypothesis or number of the counted statement or of its proof was altered. Layout and class: the article's own class is \texttt{article} with packages including graphicx, amsfonts, amsmath, amsthm, color, hyperref; the wrapper is the lane's pinned \texttt{amsart} editorial wrapper at 10pt on A4, which restates the theorem-like environments the retained text needs and adds the article's own macro definitions that the retained text needs (none), with extra wrapper packages (bm, bbm, dsfont, xspace, cancel, mathtools). The delivered numbering is the unit's own. Assets: no retained span contains a figure environment, a graphics-inclusion command, a PDF-page inclusion, a table or a TikZ picture; no image file is copied into this package, the package declares no asset dependency and no image file is redistributed with it. Licence: version arXiv:2510.24128v1 of this article is distributed under the Creative Commons Attribution 4.0 International licence, as recorded in the arXiv licence metadata for this exact version and shown on the abstract page https://arxiv.org/abs/2510.24128v1. A full rights notice is delivered at licenses/arxiv-2510.24128-CC-BY-4.0.txt. Characters: every retained excerpt is pure ASCII, so the delivered engine, XeLaTeX with its default Latin Modern text font, reports no missing character.
\begin{assumption}\label{ass}
The coefficients $b,\sigma$ and $f$ are   Lipschitz continuous  with linear growth  in  $x$, uniformly in $t$, i.e., there exists a constant $C$ such that, for any $t\in [0,T]$ and $x,y \in \mathbb R^d$, 
$$
|b(t,x)|,|\sigma(t,x)|,|f(x)| \le C(1+|x|),
$$
and
$$
|b(t,x)-b(t,y)|,|\sigma(t,x)-\sigma(t,y)|,  |f(x)-f(y)| \le C|x-y|.
$$
Moreover, $\sigma\sigma^T$ is uniformly non-degenerate, i.e., there exists a constant $c$ such that $\sigma\sigma^T(t,x) \ge cI$ for any $t\in [0,T]$ and $x \in \mathbb R^d$.  

\end{assumption}

\begin{equation}\label{extended-HJB}
	\left\{
	\begin{split}
		&\partial_t V^\lambda+\mathcal LV^\lambda+\lambda\exp(-\frac{V^\lambda+\frac{\gamma}{2}(f-g^\lambda)^2-f}{\lambda})-\gamma |\sigma \partial_x g^\lambda|^2=0,\ V^\lambda(T,x)=f(x),\\
		& \partial_t g^\lambda+\mathcal L g^\lambda-\exp(-\frac{V^\lambda+\frac{\gamma}{2}(f-g^\lambda)^2-f}{\lambda})(g^\lambda-f)=0,\ g^\lambda(T,x)=f(x).
	\end{split}
	\right.
\end{equation}

\begin{lemma}\label{append_lemma}
Let $\varphi$ be the solution of 
$$
(\partial_t+\mathcal L)\varphi+c\varphi=g, \varphi(T,x)=f,
$$
with $c,g$ being bounded continuous functions and $f$ also bounded with bounded derivatives. Then
\begin{enumerate}
    \item There exists a constant $C$ depending on the coefficients such that 
    $$
    -C(\|f^-\|_\infty+T\|g^-\|_\infty)\le \varphi \le C(\|f^+\|_\infty+T\|g^+\|_\infty),
    $$
    where $f^{\pm}$ represents the positive and negative part of $f$ respectively. If $c\le 0$, the constant $C$ can be chosen to be $1$.
    \item Assuming that $c\equiv 0$, we have 
    $$
    \|\partial_x \varphi\|_\infty \le (1+C\sqrt{T})\|\partial_x f\|_\infty+C(\sqrt{T}+T)\|g\|_\infty.
    $$
\end{enumerate}
\end{lemma}

\begin{theorem} \label{thm3.2}
 In addition to Assumption \ref{ass}, we further assume that the coefficients $b$,$\sigma$ and $f$ are uniformly bounded.  Then, for a sufficiently small $T$, \eqref{extended-HJB} admits a classical solution $(V^\lambda,g^\lambda)$.
\end{theorem}

\begin{proof}The solution is to be considered as a fixed point of a contraction mapping. For that purpose, 
denote by $\mathbb K$ the Banach space $C([0,T],C^1(\mathbb R^d))$ equipped with the norm $\|l\|_{\mathbb K}=\sup_{t,x}|l(t,x)|+\sup_{t,x}|\partial_x l(t,x)|$. For any $l \in \mathbb K$, define a mapping $F$ as $k=F(l)$ is the solution of the following system
\begin{equation}\label{decouple-HJB}
\left\{
\begin{split}
&\partial_t v+\mathcal Lv+\lambda\exp(-\frac{v+\frac{\gamma}{2}(f-l)^2-f}{\lambda})-\gamma |\sigma \partial_x l|^2=0, v(T,x)=f,\\
& \partial_tk+\mathcal L k-\exp(-\frac{v+\frac{\gamma}{2}(f-l)^2-f}{\lambda})(k-f)=0,k(T,x)=f.
\end{split}
\right.
\end{equation}
Let $B_m(0)$ be the ball in $\mathbb K$ centered at $0$ with radius $m=\|f\|_{\infty}+\|\partial_x f\|_{\infty}+1$. We are to show that, for a sufficiently small $T$, $F$ is a contraction from $B_m(0)$ into itself and, thus, admits a fixed point, which would  be the solution of \eqref{extended-HJB}. 

The proof consists of several steps:

{\it Step $1$. Well-posedness of first equation.} Choose any $N>0$, let $\xi_N$ be a smooth cutoff function such that $\xi_N(x)=x$, for $x\le N$; and $\xi_N(x)=N+1$, for $x\ge N+1$. Consider the following equation 
$$
\partial_t v+\mathcal Lv+\lambda\exp(\xi_N(-\frac{v+\frac{\gamma}{2}(f-l)^2-f}{\lambda}))-\gamma |\sigma \partial_x l|^2=0, v(T,x)=f.
$$
Noting that the third term is a bounded Lipschitz function of $v$, it admits a solution $v$. Lemma \ref{append_lemma} yields that 
$$
-C( \|f\|_{\infty}+\lambda \exp(\frac{N+1}{\lambda}))\le v\le C(Tm^2+\|f\|_{\infty}).
$$
Let us give a refined lower bound estimation independent of $N$, which implies that $v$ solves the first equation in \eqref{decouple-HJB} when $N$ is sufficiently large.
Denote by $\psi(x)=\sqrt{1+|x|^2}$.  One can compute that 
$$
D_x \psi=\frac{x}{\sqrt{1+|x|^2}}\text{ and }  D^2_x \psi= \frac{1}{\sqrt{1+|x|^2}}I -\frac{1}{\left(1+|x|^2\right)^{\frac12}} x\otimes x,
$$
and, thus, $\mathcal L \psi \le C$. Since $v$ is a bounded function, it holds that, for any $\varepsilon>0$, $v+\varepsilon\psi$ attains a minimum at some point $(t^*,x^*)$. If $t^*=T$, then the terminal condition implies that 
$$v\ge -\|f\|_\infty-\varepsilon \psi.
$$
If $t^*<T$, it holds that, at $(t^*,x^*)$,
$$
D^2_x v \ge -\varepsilon D^2_x \psi, D_x v=-\varepsilon D_x \psi,  \text{ and } \partial_t v \ge 0. 
$$
This yields that  $(\partial_t +\mathcal L)v(t^*,x^*) \ge-\varepsilon \mathcal L\psi \ge -C\varepsilon$. On the other hand, we have  
$$
(\partial_t +\mathcal L)v=\gamma |\sigma D_x l|^2-\lambda\exp(\xi_N(-\frac{v+\frac{\gamma}{2}(f-l)^2-f}{\lambda})).
$$
Combining these two inequalities, we get that 
$$
\lambda\exp(\xi_N(-\frac{v+\frac{\gamma}{2}(f-l)^2-f}{\lambda})) \le \gamma |\sigma D_x l|^2 +C\varepsilon,
$$
or, equivalently,
$$
\xi_N(\frac{f-v-\frac{\gamma}{2}(f-l)^2}{\lambda})\ \le \log \frac{\gamma|\sigma \partial_x l|^2+C\varepsilon}{\lambda}.
$$
This implies that,  at $(r^*,x^*)$,
$$
v \ge f-\lambda  \log \frac{\gamma|\sigma \partial_x l|^2+C\varepsilon}{\lambda}.
$$
Thus, we have a lower bound estimation
$$
v\ge -\|f\|_\infty +\lambda \log \frac{\gamma Cm^2+C\varepsilon}{\lambda}-\varepsilon \psi.
$$
Letting $\varepsilon$ go to $0$, we finally get that 
$$
v\ge -\|f\|_\infty +\lambda \log \frac{\gamma Cm^2}{\lambda}.
$$

{\it Step 2. Bound and derivative estimations for $k$.} Note that the second estimation of \eqref{decouple-HJB} is just a linear equation of $k$. The well-posedness is straight-forward and one obtains the following bound estimation from Lemma \ref{append_lemma} 
$$
\|k\|_{\infty} \le \|f\|_{\infty}.
$$
To give estimation of the derivative, note that, from lower bound estimation for $V$ in the first step, it holds that 
$$
\exp(-\frac{v+\frac\gamma2(f-l)^2-f}{\lambda}) \le (\frac{\gamma Cm^2}{\lambda})^{\frac{C}{\lambda}(1+\|f\|^2_\infty+m^2)}.
$$
Hence, using Lemma \ref{append_lemma} again, we have that 
$$
\|\partial_x k\|_{\infty} \le (1+\sqrt{T})\|\partial_x f\|_{\infty}+C(\sqrt{T}+T)(\frac{\gamma Cm^2}{\lambda})^{\frac{C}{\lambda}(1+\|f\|^2_\infty+m^2)}. 
$$
Then, for a sufficiently small $T$, $F$ is mapping from $B_m(0)$ into itself. 

{\it Step 3. Contraction of the mapping $F$.}
For any $h_i \in B_m(0)$ with $i=1,2$, let $(v_i,k_i)$ be the related solutions of \eqref{decouple-HJB}.  Denote by $\delta v=v_1-v_2$, $\delta k=k_1-k_2$, and $\delta l=l_1-l_2$. Then, we that $(\delta v,\delta k)$ satisfy 
\begin{equation}\label{delta-HJB}
\left\{
\begin{split}
&(\partial_t+\mathcal L)\delta v- \delta e(\delta v-\frac\gamma 2(2f-l_1-l_2)\delta l)-\gamma(|\sigma \partial_x l_1|^2-|\sigma \partial_x l_2|^2)=0, \delta v(T,x)=0,\\
& (\partial_t+\mathcal L) \delta k-\exp(-\frac{v_1+\frac{\gamma}{2}(f-l_1)^2-f}{\lambda})\delta k+(k_2-f)\delta e\frac{\delta v-\frac\gamma 2(2f-l_1-l_2)\delta l}{\lambda} =0,\delta k(T,x)=0.
\end{split}
\right.
\end{equation}
with 
$$
\delta e=\int_0^1 \exp(-\frac{v_1+\frac{\gamma}{2}(f-l_1)^2-f+s(\delta v-\frac\gamma 2(2f-l_1-l_2)\delta l)}{\lambda})ds.
$$
From the estimation in previous step, we know that  the related functions $v_i,k_i$ and $h_i$ are uniformly bounded. Thus, $\delta e$ is a bounded function. Moreover, one gets that 
$$
\|\delta e \frac \gamma2(2f-l_1-l_2)\delta l\|_\infty \le \frac\gamma 2\| \delta e(2f-l_1-l_2)\|_{\infty}\|\delta l\|_\infty \le C\|\delta l\|_\infty,
$$
and 
$$
\||\sigma D_x l_1|^2-|\sigma D_x l_2|^2 \|_\infty \le C\| \sigma D_x l_1+\sigma D_x l_2\|_\infty\| \sigma D_x l_1-\sigma D_x l_2\|_\infty \le C\| \sigma D_x l_1-\sigma D_x l_2\|_\infty.
$$
Then, using Lemma \ref{append_lemma} again, we have that there exists a constant $C$ depending on the coefficients, $m$ and $\lambda$ such that 
$$
\|\delta v\|_{\infty} \le CT(\|\delta l\|_{\infty}+\|\partial_x \delta l\|_{\infty}. 
$$
Furthermore, with a similar argument, 
$$
\|\delta k\|_\infty \le CT\|(k_2-f)\delta e\frac{\delta v-\frac\gamma 2(2f-l_1-l_2)\delta l}{\lambda}\|_\infty\le CT(\|\delta v\|_\infty+\|\delta l\|_\infty).,
$$
and
\begin{equation*}
\begin{split}
\|\partial_x \delta k\|_{\infty} \le& C(\sqrt{T}+T)(\|(k_2-f)\delta e\frac{\delta v-\frac\gamma 2(2f-l_1-l_2)\delta l}{\lambda}\|_\infty+\|\exp(-\frac{v_1+\frac{\gamma}{2}(f-l_1)^2-f}{\lambda})\delta k\|_\infty)\\
\le&C(\sqrt{T}+T)(\|\delta v\|_\infty+\|\delta l\|_\infty+\|\delta k\|_\infty).
\end{split}
\end{equation*}
Combining these estimations, we  see that $F$ is a contraction for a sufficiently small $T$.
\end{proof}

\end{document}
